\chapter*{ABSTRACT}
\thispagestyle{plain}

\vspace{0.5cm}

Blind and Low Vision (BLV) individuals face persistent challenges in recognizing visual information and interacting with physical objects due to the absence of visual cues. While existing assistive technologies provide partial support, most solutions focus on isolated tasks such as text recognition or simple object identification and lack an integrated, real-time interaction framework. This project presents Echora, an AI-powered sensory substitution system designed to enhance independence, accessibility, and task confidence for BLV users through multimodal perception and feedback.


Echora combines wearable sensing hardware with PC-based artificial intelligence to translate visual and semantic information into intuitive audio and haptic cues. The system integrates depth sensing and computer vision to support precise hand guidance, object recognition, text reading, and interaction assistance. By employing real-time multimodal sensor fusion and adaptive mode switching, Echora delivers context-aware feedback while minimizing cognitive load on the user.

All high-level processing and AI inference are performed locally on the PC, ensuring low latency, privacy preservation, and independence from cloud connectivity. The wearable components are limited to sensing and feedback, enabling a lightweight, modular, and scalable system architecture. Key features include audio feedback, haptic prompts, object and text recognition, and hand-guided interaction.

The proposed system is designed to support object-level interaction and information access rather than replace existing assistive tools. By unifying recognition, interaction, and perception within a single adaptive framework, Echora demonstrates the potential of AI-driven sensory substitution systems to significantly improve autonomy and quality of life for Blind and Low Vision users.
\clearpage
\ifPDFTeX
\chapter*{ARABIC ABSTRACT}
\thispagestyle{plain}

\vspace{0.5cm}

\textit{The Arabic abstract is rendered only when the document is compiled with XeLaTeX (Arabic text cannot be typeset by pdfLaTeX). To display it, set the compiler to XeLaTeX --- on Overleaf via Menu~$\rightarrow$ Compiler~$\rightarrow$ XeLaTeX, or locally by running \texttt{xelatex} instead of \texttt{pdflatex}.}
\else
\chapter*{\textarabic{الملخّص العربي}}
\thispagestyle{plain}

\vspace{0.5cm}

\begin{arabic}
\setlength{\parindent}{0pt}

يواجه المكفوفون وضعاف البصر تحدياتٍ مستمرة في التعرّف على المعلومات البصرية والتفاعل مع الأشياء المادية من حولهم بسبب غياب المؤشرات البصرية. وعلى الرغم من أن التقنيات المساعِدة الحالية تقدّم دعماً جزئياً، فإن معظمها يركّز على مهام منفصلة مثل قراءة النصوص أو التعرّف البسيط على الأشياء، ويفتقر إلى إطارٍ متكاملٍ للتفاعل في الوقت الحقيقي. يقدّم هذا المشروع نظام إيكورا، وهو نظام إحلالٍ حسّيٍّ مدعومٌ بالذكاء الاصطناعي صُمِّم لتعزيز استقلالية المستخدمين المكفوفين وضعاف البصر، وتيسير وصولهم إلى المعلومات، وزيادة ثقتهم في أداء المهام، من خلال الإدراك والتغذية الراجعة متعدّدَي الوسائط.

يجمع نظام إيكورا بين عتادِ استشعارٍ قابلٍ للارتداء وذكاءٍ اصطناعيٍّ يعمل على الحاسوب لتحويل المعلومات البصرية والدلالية إلى إشاراتٍ صوتيةٍ ولمسيةٍ سهلة الفهم. ويدمج النظام استشعار العمق والرؤية الحاسوبية لدعم توجيه اليد بدقّة، والتعرّف على الأشياء، وقراءة النصوص، والمساعدة في التفاعل. ومن خلال الدمج اللحظي لبيانات المستشعرات متعدّدة الوسائط والتبديل التكيّفي بين أوضاع التشغيل، يوفّر النظام تغذيةً راجعةً واعيةً بالسياق مع تقليل العبء الذهني على المستخدم.

تُجرى جميع عمليات المعالجة عالية المستوى واستدلال الذكاء الاصطناعي محلياً على الحاسوب، بما يضمن زمن استجابةٍ منخفضاً، ويحافظ على الخصوصية، ويُغني عن الاعتماد على الاتصال السحابي. وتقتصر المكوّنات القابلة للارتداء على الاستشعار وتقديم التغذية الراجعة، مما يتيح بنيةً خفيفة الوزن ومعياريةً وقابلةً للتوسّع. وتشمل أبرز الميزات التغذية الصوتية، والتنبيهات اللمسية، والتعرّف على الأشياء والنصوص، والتفاعل الموجَّه باليد.

صُمِّم النظام المقترح لدعم التفاعل مع الأشياء والوصول إلى المعلومات، لا ليحلّ محلّ الأدوات المساعِدة القائمة. ومن خلال توحيد التعرّف والتفاعل والإدراك ضمن إطارٍ تكيّفيٍّ واحد، يُبرز إيكورا الإمكانات الكبيرة لأنظمة الإحلال الحسّي المعتمدة على الذكاء الاصطناعي في تحسين استقلالية المكفوفين وضعاف البصر وجودة حياتهم بصورةٍ ملحوظة.

\end{arabic}
\fi
